% ============================================================
% Sección 1: Definición de agente
% ============================================================
\section{¿Qué es un agente?}

\begin{frame}{Una pregunta antes de empezar}
  \centering
  \Large
  ¿Qué tienen en común un termostato, \\
  un jugador de ajedrez artificial \\
  y ChatGPT usando un navegador?
  \vfill
  \normalsize
  \pause
  Los tres pueden describirse como \textbf{agentes}.
\end{frame}

\begin{frame}{Definición de agente}
  \definicion{Agente (Russell \& Norvig)}{
    Un agente es cualquier entidad que \textbf{percibe} su entorno a través
    de sensores y \textbf{actúa} sobre ese entorno a través de actuadores,
    con el fin de alcanzar un objetivo.
  }
  \vfill
  \begin{itemize}
    \item No requiere ``inteligencia''\, para ser agente: un termostato ya califica.
    \item Lo que lo hace \emph{interesante} es qué tan racional es su
      comportamiento frente al objetivo.
    \item Un \textbf{agente racional} elige, en cada instante, la acción que
      maximiza su desempeño esperado dado lo que ha percibido.
  \end{itemize}
\end{frame}

\begin{frame}{El ciclo agente–entorno}
  \centering
  \begin{tikzpicture}[node distance=2.2cm]
    \node[cajaAgente, minimum width=3.2cm, minimum height=1.6cm] (agente) {\textbf{AGENTE}\\[2pt]\footnotesize función/política};
    \node[cajaEntorno, minimum width=3.2cm, minimum height=1.6cm, right=4cm of agente] (entorno) {\textbf{ENTORNO}\\[2pt]\footnotesize mundo, tarea, usuario};

    \draw[flecha] (agente.30) -- ++(1.9,0) node[midway, above, etiquetaFlecha]{acción $a_t$} -- (entorno.150);
    \draw[flecha] (entorno.210) -- ++(-1.9,0) node[midway, below, etiquetaFlecha]{percepción $s_{t+1}$, recompensa $r_t$} -- (agente.-30);
  \end{tikzpicture}

  \vfill
  \begin{itemize}
    \item Sensores $\to$ estado percibido.
    \item Programa/política del agente $\to$ decide la acción.
    \item Actuadores $\to$ ejecutan la acción sobre el entorno.
    \item El entorno cambia, y el ciclo se repite.
  \end{itemize}
\end{frame}

\begin{frame}{Tipos clásicos de agentes (según su arquitectura)}
  \small
  \begin{itemize}
    \item \textbf{Reactivo simple}: regla condición-acción, sin memoria (termostato).
    \item \textbf{Basado en modelo}: mantiene un estado interno del mundo.
    \item \textbf{Basado en objetivos}: elige acciones que acercan a una meta.
    \item \textbf{Basado en utilidad}: pondera múltiples objetivos en conflicto.
    \item \textbf{Que aprende}: mejora su política con la experiencia (RL).
  \end{itemize}
  \vfill
  \begin{alertblock}{Adelanto}
    Un \textbf{agente LLM} añade una capa nueva: usa lenguaje natural como
    espacio de razonamiento interno, no solo como entrada/salida.
  \end{alertblock}
\end{frame}
